% !TeX root = ../main.tex
\section{Wstęp}
\label{sec:introduction}

% ------------------------------------------------------------------------
% Cel i zakres pracy:
%   - implementacja toru radaru pasywnego (PBR) opartego na estymacji transmitancji kanału (CFR/CIR)
%   - analiza porównawcza dwóch generacji Wi-Fi: IEEE 802.11a (Legacy) oraz IEEE 802.11ax (Wi-Fi 6)
%   - ewaluacja w dwóch modelach propagacyjnych: syntetycznym oraz znormalizowanym TGax Model-B
% Motywacja:
%   - wszechobecność sygnałów Wi-Fi jako darmowych iluminatorów okazji
%   - przejście ze starszego 802.11a do powszechnego obecnie 802.11ax i implikacje radarowe
%   - konieczność weryfikacji algorytmów w realistycznym clutterze wielodrogowym (TGax)
% Teza i pytania badawcze:
%   - wpływ 4-krotnie dłuższego symbolu OFDM w 802.11ax na prędkość jednoznaczną (aliasing Dopplera)
%   - wpływ poszerzonej luki nośnych DC (3 podnośne) i jej interpolacji na wyciek widmowy
%   - odporność filtracji MTI i algorytmu Coherent CLEAN na gęsty clutter wielodrogowy TGax
% Układ pracy:
%   - krótki opis zawartości poszczególnych rozdziałów (1-6)
% ------------------------------------------------------------------------

\subsection{Cel i zakres pracy}
\label{subsec:intro_scope}

\subsection{Motywacja}
\label{subsec:intro_motivation}

\subsection{Teza i pytania badawcze}
\label{subsec:intro_thesis}

\subsection{Układ pracy}
\label{subsec:intro_structure}